\documentclass[tikz,border=6pt]{standalone}
\usepackage{ctex}
\usepackage{amsmath}
\usetikzlibrary{calc, arrows.meta}

\begin{document}
\begin{tikzpicture}[scale=1.2, thick]

% 椭圆焦点弦示意图
% 椭圆 x^2/a^2 + y^2/b^2 = 1，a=2.5, b=1.5, c=sqrt(a^2-b^2)=sqrt(4)=2
% 焦点 F1(-2,0), F2(2,0)

\def\a{2.5}
\def\b{1.5}
\def\c{2.0}  % c = sqrt(6.25 - 2.25) = sqrt(4) = 2

% 坐标轴
\draw[->, gray, thin] (-3.2, 0) -- (3.2, 0) node[right, font=\footnotesize] {$x$};
\draw[->, gray, thin] (0, -2.0) -- (0, 2.0) node[above, font=\footnotesize] {$y$};
\node[font=\footnotesize, below right] at (0,0) {$O$};

% 椭圆
\draw[black, thick] (0,0) ellipse [x radius=\a, y radius=\b];

% 焦点
\coordinate (F1) at (-\c, 0);
\coordinate (F2) at (\c, 0);
\fill[black] (F1) circle [radius=1.8pt];
\fill[black] (F2) circle [radius=1.8pt];
\node[font=\footnotesize, below] at (F1) {$F_1(-c,0)$};
\node[font=\footnotesize, below] at (F2) {$F_2(c,0)$};

% 焦点弦：过 F2 的斜线，设斜率 k=0.8
% 直线 y = 0.8(x - 2)，代入椭圆方程
% 0.8^2(x-2)^2/2.25 + x^2/6.25 = 1
% 0.64(x-2)^2/2.25 + x^2/6.25 = 1
% 0.2844(x-2)^2 + 0.16 x^2 = 1
% 0.2844 x^2 - 1.1378 x + 1.1378 + 0.16 x^2 = 1
% 0.4444 x^2 - 1.1378 x + 0.1378 = 0
% x = (1.1378 ± sqrt(1.294 - 0.2449)) / 0.8889
% x = (1.1378 ± sqrt(1.049)) / 0.8889
% sqrt(1.049) ≈ 1.0242
% x1 = (1.1378 + 1.0242)/0.8889 ≈ 2.431
% x2 = (1.1378 - 1.0242)/0.8889 ≈ 0.1278
% y1 = 0.8*(2.431-2) = 0.345
% y2 = 0.8*(0.1278-2) = -1.498

\coordinate (A) at (2.431, 0.345);
\coordinate (B) at (0.1278, -1.498);

% 弦 AB（过焦点 F2）
\draw[red, thick] (A) -- (B);
\fill[black] (A) circle [radius=1.8pt];
\fill[black] (B) circle [radius=1.8pt];
\node[font=\footnotesize, above right] at (A) {$A(x_1,y_1)$};
\node[font=\footnotesize, below left] at (B) {$B(x_2,y_2)$};

% 焦半径 F2A, F2B（灰色虚线）
\draw[gray, dashed] (F2) -- (A);
\draw[gray, dashed] (F2) -- (B);

% 焦半径 F1A, F1B
\draw[gray, dashed] (F1) -- (A);
\draw[gray, dashed] (F1) -- (B);

% 标注 |AF2|
\node[red!70, font=\footnotesize] at (2.4, -0.4) {$|AF_2|$};

% 弦长公式标注（框外右下）
\node[font=\small, align=left] at (0, -2.6)
  {焦点弦弦长：$|AB| = \sqrt{1+k^2}\cdot\sqrt{(x_1+x_2)^2-4x_1x_2}$};

% 顶点
\fill[gray] (\a, 0) circle [radius=1.5pt];
\fill[gray] (-\a, 0) circle [radius=1.5pt];
\fill[gray] (0, \b) circle [radius=1.5pt];
\fill[gray] (0, -\b) circle [radius=1.5pt];

\node[font=\footnotesize, below right] at (\a, 0) {$a$};
\node[font=\footnotesize, below left] at (-\a, 0) {$-a$};
\node[font=\footnotesize, right] at (0, \b) {$b$};
\node[font=\footnotesize, right] at (0, -\b) {$-b$};

\end{tikzpicture}
\end{document}
